\section{Material inputs of the novel targets}\label{sec:16.9}

Two first-principles calculations of the dielectric tensor of ZrTe$_5$ give 20.75 and 9.8 along its slow axis, and the scattering reach of the material moves by a factor 2.8 between them (\cref{we:16-band}). Both numbers enter \cref{eq:16.1-L} exactly; no approximation of the response is involved, and a measurement of the tensor would settle the question. Most targets of this chapter are limited in the same way. The loss tangent of amorphous SiO$_2$ changes seventeenfold between preparations. And the probability that a single nuclear recoil of a few hundred eV in LiF leaves a readable F center has not been measured. \Cref{tab:16-inputs} lists, for each target, the inputs its rate needs, how well they are known and how their uncertainty moves the reach.

\subsection{Measured, computed and unmeasured inputs}\label{sec:16.9-inputs}

Some inputs of \cref{tab:16-inputs} are measured constants that vary from sample to sample: the loss tangent of a glass, the width of a donor line, the work function of a sheet on its substrate, the Luttinger parameter of a nanotube. The spread between samples is then the uncertainty of the input, and a detector that characterizes its own crystal or film removes it.

Others are computed constants, such as the dielectric tensor and the half-gap of a Dirac material, the Bloch-overlap coefficients $\mathcal{U}_{j'j}(\vect{G})$ of a donor, or the ohmic coupling $\alpha_{\mathrm{ohm}}$ of a tunneling defect. Their uncertainty is the spread between calculations, or the spread of a calculation over its functional and numerical choices. Such a spread is an estimate, not a bound: two calculations need not bracket the true value.

A third group are yields and parameters that no measurement isolates. One is the four factors of \cref{eq:16.6-yield} for single recoils of tens to hundreds of eV at millikelvin temperatures (\cref{sec:16.6-yields}); another is the density, dipole moment and deformation potential of two-level systems one by one (\cref{sec:16.7-open}), on which the estimate of \cref{sec:16.7-scat} rests through an assumed deformation potential of \SI{1}{eV}. For these inputs \cref{sec:8.8} assigns no width, and a reach is conditional on the value assumed.

\begin{table}[tbp]
  \caption{Material inputs of the novel targets, with their values, their uncertainty and its effect on the reach. The truncations of the responses are in \cref{tab:16-nuisances}.}
  \label{tab:16-inputs}
  \centering\scriptsize\setlength{\tabcolsep}{3pt}
  \begin{tabular}{@{}>{\raggedright\arraybackslash}p{0.15\linewidth}>{\raggedright\arraybackslash}p{0.38\linewidth}>{\raggedright\arraybackslash}p{0.43\linewidth}@{}}
    \toprule
    Target & Inputs (values) & Uncertainty and effect on the reach \\
    \midrule
    Dirac material (ZrTe$_5$) & background tensor $\varepsilon_{\mathrm{bg},ii} = 308.4$, 20.75, 126.1 (computed; earlier 187.5, 9.8, 90.9); velocities, cutoff \SI{0.191}{eV} and half-gap $\Delta_{\mathrm{D}} = \SI{11.8}{meV}$ of the measured-structure row (computed half-gaps 11.8--\SI{17.5}{meV}; measured \SI{11.75}{meV}) & between the two tensors $\sigmae^{\mathrm{lim}}$ moves by 2.8 to 2.9 and $\zeta^{\mathrm{lim}}$ by 1.8 to 2.0 (\cref{we:16-band}); the Dirac term inside the computed tensor, 3 to \SI{8}{\percent}, moves $\sigmae^{\mathrm{lim}}$ by 0.88; the computed half-gaps move the threshold mass along the slow axis from 7.1 to \SI{10.6}{keV} \\
    Metallic nanotubes, doped graphene & $g_c$ (measured 0.2--0.7; 0.17--0.34 for $r_{\mathrm{t}} = \SI{0.7}{nm}$ in vacuum); Fermi energy; substrate & line weight $\propto N_{\mathrm{ch}}g_c$; Fermi energy and substrate set the sheet plasmon \\
    Vacuum ejection (graphene) & work function $\Phi_{\mathrm{wf}} \simeq \SI{4.3}{eV}$ (substrate-dependent); occupied orbitals and binding energies & threshold mass $2\Phi_{\mathrm{wf}}/v_{\mathrm{max}}^2 = \SI{1.22}{MeV}$ moves with $\Phi_{\mathrm{wf}}$ \\
    Quantum dots (PbS in silica) & dot gap \SI{1.2}{eV} at \SI{3.3}{nm}; $\varepsilon_{\mathrm{in}}(\omega)$ of the dots; host $\varepsilon_{\mathrm{out}}$ & depolarization factor 0.57 to 6.7 for the cases of \cref{we:16-depol}, $\zeta^{\mathrm{lim}}$ moved by 1.33 to 0.39, sign unknown over 2.5--\SI{10}{eV} \\
    Molecular crystals (anthracene) & lowest singlet \SI{3.1}{eV} and crystal structure (measured); form factors and Frenkel coupling $\mathfrak{J}$ ($\pi$-electron model); transfer integral $t_{\mathrm{CT}}$ & model corrections in \cref{tab:16-nuisances} \\
    Dopants (Si:P) & level energies (measured); line strengths (central-cell model); line width $\ge \SI{4.22}{\micro eV}$; donor density against $\nd_{\mathrm{c}} \approx \SI{3e18}{cm^{-3}}$; local-field ratio $\mathfrak{f}^{\mathrm{loc}}_{\mathrm{D}}$ & resonant rate $\propto \nd_{\mathrm{D}}/\varGamma_{\mathrm{line}}$, $\zeta^{\mathrm{lim}} = \num{2.3e-16}$ at the $2p_0$ line for \SI{e15}{cm^{-3}}; $\mathfrak{f}^{\mathrm{loc}}_{\mathrm{D}}$ close to one, unsized \\
    Color centers (LiF) & displacement threshold $E_{\mathrm{d}} = \SI{27}{eV}$ (simulation input); site phonon energy 33.5--\SI{77.2}{meV}; yields $p_{\mathrm{FP}}$, $p_{\mathrm{surv}}$, $p_q$, $f_{\mathrm{opt}}$, $p_{\mathrm{FH}}$ & lightest mass 0.113--\SI{0.255}{GeV} for $E_{\mathrm{d}} = 10$--\SI{50}{eV}; yields unmeasured, the reach scales as $1/Y_{\mathrm{CC}}$ \\
    Bond breaking (H$_2$, molecular crystals) & bond energy \SI{4.75}{eV} and $\omega_{\mathrm{vib}} = \SI{0.546}{eV}$ (measured); cage escape $p_{\mathrm{esc}}$; yields of repulsive electronic states & $p_{\mathrm{esc}}$ and electronic yields unmeasured \\
    Glasses (a-SiO$_2$) & $\tan\delta_0$; $\varepsilon_r = 3.9$ (assumed); acoustic tunneling strength $C_{\mathrm{l}}$ (measured); $\bar P$, $\gamma_{\mathrm{l}}$ ($\SI{1}{eV}$ assumed), dipole moment & $\tan\delta_0 = 3\times10^{-4}$ (thermal oxide) to $5\times10^{-3}$ (deposited oxide): absorption rate $\times16.7$, $\zeta^{\mathrm{lim}}$ $\times4.1$; scattering host term fixed by $C_{\mathrm{l}}$ alone, core term needs $\bar P$ and $\gamma_{\mathrm{l}}$ \\
    Two-level systems at millikelvin & $\alpha_{\mathrm{ohm}} = 4K_{\mathrm{b}}$ (Born), $K_{\mathrm{b}} = 0.07$, 0.11 fitted (aluminum); $\bar P$ of vitreous silica $2.0\times10^{44}$ (heat release) or $8.0\times10^{44}\,\si{J^{-1}.m^{-3}}$ (specific heat) & heat release fourfold between the two $\bar P$; first-principles $\alpha_{\mathrm{ohm}}$ needs the scattering phases of both configurations \\
    \bottomrule
  \end{tabular}
\end{table}

\subsection{Theory-limited and data-limited targets}\label{sec:16.9-limits}

A target is limited by theory when its response has no controlled approximation at the values its inputs take, so that better inputs would not sharpen the rate. Dirac cones at strong coupling (\cref{sec:16.1-validity}), donors near the Mott density (\cref{sec:16.5-band}) and interacting two-level systems (\cref{sec:16.8-open}) are of this kind.

A second group is limited by a calculation whose formalism is known. The electron--hole ladder near the edge of a Dirac material changes the rate of ZrTe$_5$ by the Sommerfeld factor, 1.5 or more, at \SI{20}{keV} and by less above \SI{100}{keV} (\cref{sec:16.1-validity}). So do the final state of an ejected electron, the periodic Bethe--Salpeter response of a molecular crystal at finite $\vect{q}$, the excitons of a semiconducting tube and the valley-orbit coefficients of a donor. The response of a layered stack beyond its effective medium and the ohmic coupling $\alpha_{\mathrm{ohm}}$ from first principles complete the group. In each case the next term has a definite form and is computable, but no computation of it exists for the material.

The remaining targets are limited by data. Their inputs are properties of the sample that a measurement fixes: the dielectric tensor and the gap of ZrTe$_5$, the dielectric function $\varepsilon_{\mathrm{in}}(\omega)$ of real quantum dots, the oscillator strengths of the donor lines, the displacement threshold of LiF as a function of direction, the four yields of a color center, and the parameters of the two-level systems of a given film; \cref{sec:16.6-yields} names the measurements of the yields. For these inputs the response formula is in hand, and the uncertainty of the rate is that of the sample.

\subsection{Reach bands from input uncertainties}\label{sec:16.9-band}

An input $x$ on which the rate depends as a power, $\mathcal{R} \propto x^{-p}$ over the energies the rate samples, is a normalization uncertainty in the sense of \cref{sec:8.8}. By \cref{eq:1.5-limscaling} it moves a scattering reach as $\sigmae^{\mathrm{lim}} \propto x^{p}$, so a range $x_1 < x < x_2$ becomes a band of constant width $(x_2/x_1)^p$ in $\sigmae^{\mathrm{lim}}$ wherever $p$ does not change with the mass. An absorption reach moves as $\mathcal{R}^{-1/2}$, and its band is half as wide on a logarithmic scale. An input that sets a threshold acts instead through the low-mass wall, $\delta\ln m_{\mathrm{min}} = \delta\ln E_{\mathrm{th}}$ by \cref{eq:1.3-mmin}, and near the wall its band widens without limit (\cref{eq:8.8-shift}). An input without a size, such as the yield of a color center, gives no band: the reach is drawn for an assumed value and moves as its inverse.

For a Dirac material the power follows from \cref{eq:16.1-L}. At fixed Dirac parameters,
\begin{equation}
  \frac{\partial\ln L}{\partial\ln\varepsilon_{\mathrm{bg}}} = -\frac{2\,\varepsilon_{\mathrm{bg}}\Real(\varepsilon_{\mathrm{bg}} + \varepsilon_{\mathrm{D}})}{|\varepsilon_{\mathrm{bg}} + \varepsilon_{\mathrm{D}}|^2} = -2\Bigl[1 - \frac{\Real\varepsilon_{\mathrm{D}}}{\varepsilon_{\mathrm{bg}}} + O\Bigl(\frac{|\varepsilon_{\mathrm{D}}|^2}{\varepsilon_{\mathrm{bg}}^2}\Bigr)\Bigr] ,
  \label{eq:16.9-power}
\end{equation}
since $\varepsilon_{\mathrm{bg}}$ is real and $\varepsilon_{\mathrm{D}}$ does not depend on it. The reach therefore moves almost as $\varepsilon_{\mathrm{bg}}^2$, less by twice the fraction of the screening that the Dirac electrons supply. For ZrTe$_5$ that fraction is 3 to \SI{8}{\percent} at long wavelength (\cref{we:16-band}).

\begin{workedexample}{The reach band of ZrTe$_5$ from its background tensor}\label{we:16-band}
Take the Dirac parameters of \cref{sec:16.1-zrte5}, the measured-structure row with $\Delta_{\mathrm{D}} = \SI{11.8}{meV}$, and the scattering reach of \cref{fig:16-zrte5}a: light mediator, orientation-averaged loss function, three events in a kg-year. The central tensor is $\varepsilon_{\mathrm{bg},ii} = 308.4$, 20.75 and 126.1~\cite[Table~2]{geilhufe2020dirac}. The earlier calculation of \cref{we:7-zrte5} gives 187.5, 9.8 and 90.9~\cite[Table~1]{hochberg2018detection}, smaller by factors of 0.61, 0.47 and 0.72. We combine it with the same velocities, gap and cutoff, which isolates the effect of the tensor; the two calculations also differ in their velocities.

\emph{Estimate.} At weak screening $L \to \Imag\varepsilon_{\mathrm{D}}/(\hat{\vect{q}}\cdot\varepsilon_{\mathrm{bg}}\cdot\hat{\vect{q}})^2$. If $\Imag\varepsilon_{\mathrm{D}}$ were the same in every direction, the reach would move by the ratio of the direction averages of $(\hat{\vect{q}}\cdot\varepsilon_{\mathrm{bg}}\cdot\hat{\vect{q}})^{-2}$, which is 0.327 for the earlier tensor against the central one. Directions within \ang{30} of the slow axis $y$, where the background is smallest, carry \SI{64}{\percent} of that average for the central tensor, so the band is set mainly by $\varepsilon_{\mathrm{bg},yy}$, $(9.8/20.75)^2 = 0.22$ along the axis itself.

\emph{Full calculation.} With \cref{eq:16.1-L} and the anisotropic $\varepsilon_{\mathrm{D}}$ of \cref{eq:D.41-aniso}, the earlier tensor gives a reach 0.344 to 0.362 times the central one between \SI{20}{keV} and \SI{1}{MeV}, and 0.31 just above threshold. At the best mass, \SI{36}{keV}, the reach lies between $\SI{1.53e-40}{cm^2}$ and $\SI{5.5e-41}{cm^2}$, a band of 2.8 that is nearly constant in mass. The full ratio lies above the estimate because the smaller tensor lets the Dirac electrons screen more: the effective coupling along $y$ rises from 0.144 to 0.305.

\emph{Power.} Scaling the whole central tensor by 1.1 moves the reach by $1.1^p$ with $p = 1.86$ to 1.94 between \SI{20}{keV} and \SI{1}{MeV}, and $p = 1.85$ at \SI{12}{keV}, as \cref{eq:16.9-power} predicts for $\Real\varepsilon_{\mathrm{D}}/\varepsilon_{\mathrm{bg}}$ of a few percent. A tensor uncertain by a common factor of two would give a band of 3.6 to 3.8.

\emph{Dirac content.} Removing the static Dirac term of \cref{eq:7.3-aniso}, 10.0, 1.67 and 8.3 for this row, 3 to \SI{8}{\percent} of the components, leaves the background 298.4, 19.08 and 117.8, and moves the reach by 0.881 to 0.885, a correction smaller than the band.

\emph{Absorption.} For dark-photon absorption the earlier tensor moves $\zeta^{\mathrm{lim}}$ by 0.49 to 0.55 between the gap and the cutoff, from $1.24\times10^{-15}$ to $6.4\times10^{-16}$ at \SI{100}{meV}. The polarization along $y$ carries most of the rate, and for it alone $\zeta^{\mathrm{lim}} \propto \varepsilon_{\mathrm{bg},yy}$ at weak screening, which would give 0.47.

\emph{Gap.} The computed half-gaps of 11.8 to \SI{17.5}{meV} put the threshold mass along $y$ between 7.1 and \SI{10.6}{keV}, with \SI{9.4}{keV} for the relaxed-cell row at \SI{15.6}{meV}; the measured \SI{11.75}{meV} gives \SI{7.1}{keV}. Close to the wall the threshold band exceeds the tensor band, since it diverges there.
\end{workedexample}

Above \SI{100}{keV}, where the region near the pair edge carries less than a fifth of the rate, the tensor band is the largest sized uncertainty of the ZrTe$_5$ reach.
